\section{相關研究}

本研究主要位於四條研究脈絡的交會處：noticed-case retrieval、THUIR 的結構化法律預訓練、檢索互動架構與排序目標，以及長文本 encoder 與動態負樣本設計。

\subsection{法律判決檢索與 THUIR@COLIEE 2023}

COLIEE Task~1 的重點不是一般的相似案例搜尋，而是要找出真正支撐 query case 判決理由的 noticed cases。THUIR 在 COLIEE 2023 \cite{li2023thuir} 以 SAILER 為核心，結合 reference sentence extraction、post-processing 與 learning-to-rank，取得當年度最佳成績。其最重要的貢獻，是提出結構感知的法律預訓練：將判決書拆成 Fact、Reasoning、Decision 三段，由較深的 encoder 編碼 Fact，再把 `[CLS]` hidden state 注入兩個較淺的 decoders，用來重建被遮罩的 Reasoning 與 Decision。最終訓練目標為三個 masked language modeling loss 的總和，其中 Fact 的遮罩比例較低，Reasoning 與 Decision 較高，目的是讓 encoder 更專注於把案件事實壓縮成可檢索的表徵 \cite{li2023thuir}。本文延續 THUIR 的問題意識，但把焦點轉向兩個更直接的問題：citation-context query 是否仍適合長文本 encoder，以及 dense retrieval 是否主要受限於截斷。

\subsection{檢索互動架構與排序目標}

在神經式檢索架構中，dual encoder 會將 query 與 candidate 分別編碼成向量，再以向量相似度計分，因此可以離線預先建立 candidate embeddings，是 dense retrieval 最常見的第一階段架構 \cite{karpukhin2020dpr}。cross encoder 則把 query 與 candidate 串接後一起輸入同一個 Transformer，讓所有 token 直接交互，通常更適合 reranking，但需要對每個 query--candidate pair 個別推論，成本顯著較高 \cite{nogueira2019bertpr}。late interaction 介於兩者之間，例如 ColBERT 先分別編碼 query 與文件 token，再以較便宜的 token-level interaction 聚合相似度，在效果與效率之間取得折衷 \cite{khattab2020colbert}。

這種差異對本研究特別重要。由於 query 與 candidate 都是長篇判決書，若採用 cross encoder，單次 reranking 在單張 RTX 4080 SUPER 16GB 上很容易因 token 數暴增而超出模型上限或造成過高的記憶體成本。因此，本文把 dual encoder 作為第一階段主體，並把 cross encoder 與 late interaction 保留為後續工作。相對地，後段排序則採用 learning-to-rank。pointwise 方法把每個 query--document pair 視為獨立樣本；pairwise 方法學習哪一篇文件應排在另一篇前面，例如 RankNet \cite{burges2005ranknet}；listwise 方法則直接對整個候選排序進行學習，更貼近真實 ranking metric \cite{cao2007listnet,liu2009ltr}。本文的 dual encoder 對比式學習較接近 listwise：雖然 query 與單一文件的相似度是逐對計算，但每一步的 InfoNCE 損失會把 1 個 positive 與 15 個 negatives 放進同一個 softmax 共同正規化，因此優化對象是整個候選集合中的相對排序，而不是彼此獨立的二元 pair。相對地，LightGBM LambdaRank 在更新形式上屬於 pairwise，因為它透過文件對之間的交換來產生梯度；但這些 lambda 的權重又直接來自整體排序對 NDCG 的影響，因此其最終優化目標帶有明顯的 listwise 性質。

\subsection{長文本 Encoder 與 ModernBERT}

ModernBERT \cite{warner2024modernbert} 是近年對 encoder-only Transformer 的重要改良模型，其設計重點不只是把 context window 拉長，而是同時兼顧記憶體效率與推論速度。它採用 RoPE、pre-norm 與 GeGLU，並移除多數 bias 項；在注意力設計上，交替使用 global attention 與 local sliding-window attention，其中每第三層使用 global attention，其餘層使用 128-token local attention；在實作層則搭配 unpadding、Flash Attention 2/3 與編譯優化，以提升長文本訓練與推論效率 \cite{warner2024modernbert}。對本文而言，ModernBERT 的價值在於它提供了一個可直接檢驗「長文本是否為法律判決 dense retrieval 主要瓶頸」的長上下文 encoder。

\subsection{動態負樣本與合法候選範圍}

Dense retrieval 的對比式學習高度依賴負樣本品質。若負樣本過於容易，模型學不到穩定的判別邊界；若負樣本來源固定，卻又與當前模型狀態脫節，訓練同樣可能失敗。本文早期依照類似 THUIR 的流程，先以 BM25 top-100 建立 hard-negative pool，再在每個步驟抽取 15 篇負樣本，但幾乎無法收斂。這顯示法律判決檢索中的 hard negatives 不能只由 lexical rank 決定，還必須隨模型狀態動態更新。

此外，法律檢索還有一個一般網頁檢索較少遇到的問題：若不限制候選範圍，某些語義相關、但年份晚於 query case 的未來判決，會被誤當成負樣本。這些文件未被標成 noticed case，未必因為它們真的不相關，而可能只是法官在時間上根本不可能引用。因此，本文把動態負樣本選取與 query-specific 年份 scope 結合，讓模型只在法律上合理的候選範圍內學習真正的負樣本。
